% latexmk -pdf -pvc blatt14.tex
\documentclass{hott-übung}

\begin{document}

\setcounter{blattnummer}{14}
\setcounter{aufgabennummer}{0}

\blatt

\aufgabe{}
Wir definieren die strikte Ordnungsrelation \(<\) auf \(\N\) per Induktion durch
\begin{align*}
  (0 < 0) &\colonequiv \leer, & (0 < n + 1) &\colonequiv \eins, \\
  (m + 1 < 0) &\colonequiv \leer, & (m + 1 < n + 1) &\colonequiv m < n.
\end{align*}
Für \(n : \N\) war die standard-endliche Menge mit \(n\) Elementen durch
\[\{0,\dots,n-1\} \equiv \sum_{k : \N} k < n\]
gegeben. Weiter sei der abhängige Typ \(\mathrm{Fin} : \N \to \mU\) per Induktion definiert durch
\begin{align*}
\mathrm{Fin}(0) \colonequiv \leer, && \mathrm{Fin}(n+1) \colonequiv \eins \sqcup \mathrm{Fin}(n).
\end{align*}
\begin{enumerate}
  \item
    Finde eine Äquivalenz \(\mathrm{Fin}(n) \simeq \{0,\dots,n-1\}\) für alle \(n : \N\).
  \item
    Zeige, dass jeder endliche Typ \(E\) das \emph{Auswahlaxiom} erfüllt:
    Für alle \(A : E \to \mU\) gilt
    \[\prod_{x:E}\|A(x)\| \to \|\prod_{x:E}A(x)\|.\]
\end{enumerate}

\aufgabe{}
\begin{enumerate}
\item Zeige, dass \(S^1\vee S^1\) ein endlicher CW-Komplex ist.
\item Berechne \(H^1(S^1\vee S^1)\) und folgere \(\forall n:\N.\neg (S^1\vee S^1\simeq S^n)\).
  {\tiny Tipp: Mayer-Vietoris.}
\end{enumerate}

\aufgabe{}
Es sei \(d : S^1 \to S^1\) die \enquote{Orientierungsüberlagerung} aus Aufgabe 4 von Blatt 7.
\begin{enumerate}
  \item Unter Verwendung des Isomorphismus \(H^1(S^1) \simeq \Z\) aus Aufgabe 1 von Blatt 13, berechne das Bild der induzierten Abbildung \(d^* : H^1(S^1) \to H^1(S^1)\).
  \item Es sei \(\mathrm{RP}^2\) der Pushout von \(S^1 \xleftarrow{d} S^1 \to \eins\).
    Berechne \(H^1(\mathrm{RP}^2)\).
\end{enumerate}


\end{document}
